\chapter{文本与声音：条件特征和音频表示}
\label{seedance-encoders}

视频 VAE 已经给出了画面表示。总图输入侧还有文字与声音，
它们并非先变成视频像素再进入主干。下面继续按各自的输入、编码结果和用途来写。
具体规格沿用\hyperref[seedance-m05]{联合音视频章节}与\hyperref[seedance-m21]{模块接口}登记的
LTX-2 原系列，本文没有据此认定 Seedance 使用相同组件。

\section{文本编码器与主干投影分别做什么}
先取提示词的 token 序列 $w_{1:L}$，经配套文本编码器 $T$，得到
\[
H_c=T(w_{1:L})\in\mathbb R^{B\times L\times d_c}.
\]
这里特征宽度是文本编码器的输出宽度，并不要求等于视频流或声音流的隐藏宽度。
连接器与投影把它变成主干能够读取的条件，随后在相应文本 attention 中生成 K/V，
视觉或音频状态生成 Q。若使用联合多模态自注意力，文本也可以产生自己的 Q；
因此两种连接方式应按实际基座区分，不能只凭“文本条件”猜接口。

本文公开替代基座保留 Gemma 文本组件与配套连接器。冻结文本模型，
表示 $T$ 的参数不更新；新投影或参考实体汇聚是否更新，仍由各自参数组决定。
这两件事若混在一起，就容易把可学习支路误写成整个冻结条件。

另一个问题是角色词怎样对应参考图。正文 caption 可以继续走原生文本路径，
实体描述则先在原始 token 序列定位，汇聚后放入独立实体记忆，
与相应图像附同一个主体槽和 role。
连接器可能重组序列、插入 register，因此字符跨度不能直接当作连接器后的 token 索引。
具体绑定在\hyperref[seedance-m22]{参考身份与时间路由}中展开。

\section{声音的采样率与潜时间率}
声音首先是波形 $x^a[n]$，物理时间由 $\tau=n/f_s$ 给出。
codec 再将一段波形或配套声学表示映射到潜变量，潜时间位置有自己的速率。
两个速率分别描述原始采样与潜表示，不能把每秒样本数当成主干的 token 数。

LTX-2 原系列的接口例子中，音频输入使用 16kHz mel 表示，
codec 可保留 $B\times8\times T_a\times16$ 的布局。
进入主干前，将每个时间位置的两个特征轴整理为 $8\times16=128$ 维，
约每秒 25 个时间位置，再由输出端 vocoder 得到 24kHz 立体声。
这是一套配套输入输出约定，不表示任意 16kHz 编码器都能接相同解码器。

对时长 $D$，可以先估 $T_a\approx25D$，但实际还要遵守具体编码器的边界与长度规则。
每个潜 token 保存其覆盖区间，再与视频 VAE 的时间网格对齐。
视频帧、视频 latent、音频采样、音频 latent 的下标都不同，
跨模态 attention 应比较物理时间，而不是直接比较下标差。

\section{同一段声音在不同任务里是什么变量}
设编码结果为 $z_0^a=E_a(x^a)$。它可以有不同用途。
联合生成时，声音是未知目标，构造 $z_t^a$ 并回归相应速度；
音色参考时，它被汇聚成条件记忆，目标台词另行给定；
保留原音轨时，它是干净条件，采样器不再更新音频 latent，交付时复用原波形。

所以一个“音频已提供”的布尔变量还不够。
是否存在音轨、是否计算声音隐藏状态、是否更新声音目标、是否计算音频 loss，
分别控制不同路径。缺音轨没有目标，真实静音却仍有监督；
保持声音不更新，也不要求声音流停止给视频提供条件。
这些状态在\hyperref[seedance-m23]{音画编辑保持}中会逐项对应。

最后，声音隐藏状态可能读取当前视频。输入音轨固定以后，codec 特征可复用，
深层 Transformer 声音状态仍可能逐步变化。因此缓存结论要沿计算图继续追到依赖，
不能只停在源音频文件没有改变这一层。
